\begin{lemma}[A terminal common-missing color]
Fix a finite loopless multigraph with individual edge identities, a
palette of $k$ colors, and a distinguished edge $e_0$.  Let
$(T,\varphi_n)$ be a positive-rung extended Tashkinov tree constructed
from an arbitrary certified input coloring, let $\mathcal B$ be a
hierarchy base under $\varphi_n$, and let $C$ be a hierarchy
counterexample over that base.  Assume that $\varphi_n$ is proper and
leaves exactly $e_0$ uncolored.

Assume, uniformly for every positive-rung maximum ETT on the same graph,
palette, and distinguished edge, with arbitrary certified input and
generating colorings and with smaller rung number, all six lower-rung
assertions of the technical theorem: elementarity; terminal
interchangeability; at every PE slot, unconditional stability before the
support, stability on the preceding ladder tree when the first-appearance
index is greater than one, and disjointness of stored pairs at distinct PE
slots with distinct supports and equal first-appearance index; the terminal
PE component assertion; the defective-endpoint order assertion; and the
mod-coloring assertion.  Assume also that every
same-rung maximum ETT with a good hierarchy having fewer dividers than
$C$ is elementary.

Assume that $C$ is lexicographically minimal among hierarchy
counterexamples whose ETTs have the same rung number as $T$ and whose
hierarchies have the same divider count, allowing the certified input
coloring, generating coloring, ETT, base, and hierarchy to vary.  Finally
assume
\[
 2n+11\leq |\overline\varphi_n(T_n)|,
\]
where $n$ is the rung number and $T_n$ is the terminal ladder tree.

Write $H=C.\mathit{hierarchy}$ and
\[
 A=\noderef{HierarchyStarredStage}
      (H,\operatorname{Fin.last}(H.\mathit{dividerCount})),
 \qquad \ell=|T|-|A|.
\]
Let
\[
 B=\noderef{TreeSequencePrefix}
       \bigl(T,\min\{|A|+(\ell-1),|T|\}\bigr)
\]
be the exact ordered penultimate full-tail prefix, and let $y_\ell$ be
the final displayed vertex of $T$.  Then there is an actual color
$\alpha$ such that
\[
 \alpha\in\overline\varphi_n(B)
       \cap\overline\varphi_n(y_\ell).                         \tag{7.1}
\]
The prefix $B$ inherits all displayed edge identities from $T$, including
the identities of parallel edges.
\end{lemma}
\begin{proof}
Apply \noderef{HierarchyCounterexampleTailAtLeastTwo} to the complete
ambient package in the statement.  It gives
\[
                              2\leq\ell.                       \tag{1}
\]
Apply \noderef{HierarchyCounterexamplePenultimateElementary} to the same
package.  Its conclusion is exactly
\[
             V(B)\text{ is elementary under }\varphi_n.        \tag{2}
\]

Since $\ell=|T|-|A|$ and (1) implies $|A|\leq |T|$, elementary arithmetic
gives
\[
 |A|+\ell=|T|,
 \qquad
 \min\{|A|+\ell-1,|T|\}=|T|-1.                                \tag{3}
\]
Consequently \noderef{TreeSequencePrefix} says that $B$ consists of the
first $|T|-1$ exact edges and the first $|T|$ vertices of $T$, entry by
entry.  We verify the resulting vertex-set decomposition directly.
By \noderef{TreeSequenceVertexSet}, $V(B)$ is the image of the vertex
map of $B$ over all its vertex indices.  The defining restricted vertex
map in \noderef{TreeSequencePrefix} sends an index
$i\in\operatorname{Fin}(|T|)$ to the vertex of $T$ with the same
numerical index $\operatorname{val}(i)<|T|$.  On the other hand, the
index of the final vertex $y_\ell$ in $T$ is
$\operatorname{Fin.last}(|T|)$, whose numerical value is $|T|$.  If
$y_\ell$ belonged to $V(B)$, the image definition
would therefore give an $i$ with $\operatorname{val}(i)<|T|$ for which
the two corresponding values of the vertex map of $T$ agree.  The
vertex injectivity field of
\noderef{TreeSequence} would make their indices equal, contradicting
$\operatorname{val}(i)<|T|$.  Hence
\[
                         y_\ell\notin V(B).                    \tag{4}
\]

Conversely, let $u\in V(T)$.  Again unfolding
\noderef{TreeSequenceVertexSet}, write $u=T.\mathit{vertex}(j)$ for
$j\in\operatorname{Fin}(|T|+1)$.  If $\operatorname{val}(j)<|T|$, the
index with the same numerical value in $\operatorname{Fin}(|T|)$ maps
under the restricted vertex map of \noderef{TreeSequencePrefix} to $u$,
so $u\in V(B)$.  If
$\operatorname{val}(j)\not<|T|$, then
$\operatorname{val}(j)<|T|+1$ forces $\operatorname{val}(j)=|T|$, hence
$j=\operatorname{Fin.last}(|T|)$ and $u=y_\ell$.  We have thus proved
the exact alternatives
\[
 u\in V(T)\quad\Longrightarrow\quad
 u\in V(B)\ \hbox{or}\ u=y_\ell.                              \tag{5}
\]
The edge map in \noderef{TreeSequencePrefix} is restricted in the same
entrywise fashion, which is why no parallel edge identities are
identified in this argument.

Suppose that (7.1) were false.  We show that $V(T)$ would then be
elementary.  Let $u,v$ be distinct vertices of $T$.  By the preceding
decomposition (5), each belongs either to $V(B)$ or is $y_\ell$.  They cannot
both be $y_\ell$.  If both belong to $V(B)$, their missing-color sets are
disjoint by (2) and \noderef{ColoredVertexSetElementary}.  In the remaining
case, after interchanging $u$ and $v$ if necessary, we have
$u\in V(B)$ and $v=y_\ell$: indeed, (4) ensures that the vertex assigned
to $V(B)$ is not also $y_\ell$.  If a color $\alpha$ were missing at both,
then the union definition \noderef{TreeMissingColors} would put
$\alpha$ in $\overline\varphi_n(B)$, contrary to the assumed failure of
(7.1).  Thus the two missing-color sets are disjoint in this case as well.

We have proved pairwise disjointness of the missing-color sets at every
two distinct vertices of $T$.  By
\noderef{ColoredVertexSetElementary}, this says that $V(T)$ is elementary
under $\varphi_n$, contradicting the `notElementary' field of
\noderef{HierarchyCounterexample}.  Hence (7.1) holds, and choosing a
member of its nonempty finite intersection supplies the asserted color
$\alpha$.
\end{proof}
